\documentstyle[%


righttag%


,11pt%


,amssymb%


,russian%               remove in Eng.


,izvru%                 Set izven% in Eng.


% ,izvru-ks             for brief communications


]{amsart}





%       Math definitions





\newcommand{\la}{\langle}


\newcommand{\ra}{\rangle}


\newcommand{\tts}[1]{}





%\hoffset=-15mm                  For Eng.





\setcounter{page}{15}


\god{1997}


\nomer{11~(426)} %                Remove in Eng.


%\nomer{Vol.\,41, No.\,11}        Set in Eng.


% \str{75--81}                    Set in Eng.


\udk{511.56}


\author{\it А.А.\,Вьялицин}


\title{Исследование систем линейных диофантовых\\


неравенств методами теории групп}





%\thanks{Работа выполнена при поддержке РФФИ, грант \No. 94-01-00183.}


\date{}


% \pagestyle{myheadings}         Set in Eng.





\pagestyle{plain} %              Remove in Eng.


\begin{document}





% \thispagestyle{plain}          Set in Eng.





\maketitle





% Body text





Рассматривается система линейных диофантовых неравенств


\begin{align}


&\la\bold a^i,\bold x\ra\le b_i,\quad i=1,\dotsc,l;\\


&\bold x\in Z^n,


\end{align}


где $Z^n$ --- множество целочисленных векторов арифметического


пространства $R^n$, $\la\bold a^i,\bold x\ra$ --- скалярное произведение


векторов $\bold a^i$, $\bold x$.





Предполагается, что множество решений системы неравенств (1) непусто и


ограничено. Условие (2) --- это инвариант группы движений целочисленной


решетки $Z^n$. Поэтому есть смысл применить методы теории групп при


исследовании системы (1), (2).





Пусть $\Lambda^0$ --- группа подстановок на множестве


$\{1,\dotsc,n\}$. Обозначим через


$\bold x\lambda=(x_{\lambda(1)},\dotsc,\linebreak


x_{\lambda(n)})$


результат действия подстановки


$\lambda=(\lambda(1),\dotsc,\lambda(n))$ на вектор


$\bold x=(x_1,\dotsc,x_n)$. В этом смысле


группа


$\Lambda^0$ --- это конечная подгруппа группы движений целочисленной


решетки и одновременно конечная подгруппа ортогональной группы.





Вначале симметризуем систему (1). Положим


$C=\{\bold a^i\lambda\mid\forall\lambda\in\Lambda^0,\;


i=1,\dotsc,l\}=\{\bold c^1,\dotsc,\bold c^m\}$;


$\bold c^i=\bold a^i$, $i=1,\dotsc,l$. Определим


числа $d_i$ следующим образом. Если


$i=1,\dotsc,l$, то $d_i=b_i$; если $i=l+1,\dotsc,m$, то


$d_i=\max\{\la\bold c^i,\bold x\ra\mid\la\bold a^k,\bold x\ra\le b_k,\;


k=1,\dotsc,l\}$.


Система


\begin{equation}


\la\bold c^i,\bold x\ra\le d_i,\quad i=1,\dotsc,m,


\end{equation}


эквивалентна системе (1). По определению множество $C$ симметрично по


группе $\Lambda^0$, т.е.


$C\Lambda^0=\{\bold c^i\lambda\mid\forall\lambda\in\Lambda^0,\;


i=1,\dotsc,m\}=C$. Будем


говорить, что системы (3)


и (2), (3) симметричны по группе $\Lambda^0$.


Если система симметрична по


группе $\Lambda^0$, то она симметрична по любой ее подгруппе.





Действие группы $\Lambda^0$


переставляет векторы $\bold c^i$. Подстановке


$\lambda\in\Lambda^0$ соответствует


подстановка $\pi=(\pi(1),\dotsc,\pi(m))$ такая, что


$\bold c^i\lambda=\bold c^{\pi(i)}$. Множество таких


подстановок образует


группу $\Pi^0$ подстановок на множестве $\{1,\dotsc,m\}$, порожденную


группой $\Lambda^0$ и множеством $C$.





\begin{lem}


Группы $\Lambda^0$, $\Pi^0$ изоморфны.


\end{lem}





\begin{pf}


Так как множество решений системы (3) непусто и ограничено, то множество


$C$ содержит $n$ линейно независимых векторов. Пусть это будут векторы


$\bold c^1,\dotsc,\bold c^n$. Отображение


$\Lambda^0$ на $\Pi^0$ --- это гомоморфизм.


Пусть $\lambda$ принадлежит ядру


гомоморфизма. Представляя любой вектор $\bold x\in R^n$ линейной


комбинацией векторов $\bold c^1,\dotsc,\bold c^n$, получим


$$


\bold x\lambda=\biggl(\,\sum^n_{i=1}\alpha_i\bold c^i\biggr)\lambda


=\sum^n_{i=1}\alpha_i(\bold c^i\lambda)=


\sum^n_{i=1}\alpha_i\bold c^i=\bold x.


$$


Следовательно, ядро гомоморфизма содержит лишь единицу группы


$\Lambda^0$. Поэтому


группы $\Lambda^0$, $\Pi^0$ изоморфны.


\end{pf}





Пусть $\Lambda\le\Lambda^0$; $\Pi$


--- группа подстановок, порожденная подгруппой $\Lambda$ и


множеством $C$. Обозначим:


$P_k$, $k=1,\dotsc,p$, --- орбиты группы $\Lambda$ в множестве


$\{1,\dotsc,n\}$; $Q_r$,


$r=1,\dotsc,q$, --- орбиты группы $\Pi$ в множестве


$\{1,\dotsc,m\}$; $j(k)=\min\{j\mid j\in P_k\}$;


$f^r_k=\sum\limits_{i\in Q_r}c^i_{j(k)}$;


$\bold f^r=(f^r_1,\dotsc,f^r_p)$; $g_r=\sum\limits_{i\in Q_r}d_i$.





Запишем следующую систему линейных диофантовых неравенств:


\begin{align}


&\la\bold f^r,\bold y\ra\le g_r\quad r=1,\dotsc,q;\\


&\bold y\in Z^p.


\end{align}





\begin{lem}


Множество решений системы {\em (4)} непусто и ограничено.


\end{lem}





\begin{pf}


Введем в рассмотрение систему


\begin{align}


&|Q_r|\la\bold c^i,\bold x\ra\le g_r,\quad i\in Q_r,\quad


r=1,\dotsc,q;\\


&x_j=x_s \; \forall j,s\in P_k;\quad k=1,\dotsc,p.


\end{align}


Условия (7) определяют в $R^n$ множество неподвижных по группе


$\Lambda$


векторов: $\bold x\lambda=\bold x$ $\forall\lambda\in\Lambda$, если и


только если выполняются условия (7).





Пусть $X_1$, $X_2$, $X_3$, $X_4$ --- множества решений систем (3), (4),


(6) и (6),


(7) соответственно. Множество $X_3$ ограничено, поскольку множество


$X_1$


непусто и ограничено (\cite{Er}, с.114).





Пусть $\bold v\in X_1$. Докажем, что вектор


$$


\bold u=\frac{1}{|\Lambda|}\sum_{\lambda\in\Lambda}\bold v\lambda


$$


принадлежит множеству $X_4$. Так как


$\Lambda\lambda=\Lambda$ $\;\forall\lambda\in\Lambda$, то


$\bold u\lambda=\bold u$ $\;\forall\lambda\in\Lambda$.


Следовательно,


вектор $\bold u$ удовлетворяет условиям (7). Имеем


$\la\bold c^i,\bold u\ra=\la\bold c^i,\bold u\lambda\ra=


\la\bold c^i\lambda^{-1},\bold u\ra=\la\bold c^s,\bold u\ra$


$\;\forall i\in Q_r$, $\forall\lambda\in\Lambda$, где


$s\in Q_r$, $\lambda^{-1}$ --- подстановка, обратная подстановке


$\lambda$. Поэтому $\forall s\in Q_r$ получим


\begin{equation}


|Q_r|\la\bold c^s,\bold u\ra=\sum_{i\in Q_r}\Bigl\la\bold c^i,


\frac{1}{|\Lambda|}\sum_{\lambda\in\Lambda}\bold v\lambda\Bigr\ra=


\frac{1}{|\Lambda|}\sum_{\lambda\in\Lambda}\;


\sum_{i\in Q_r}\la\bold c^i\lambda^{-1},\bold v\ra \le


\frac{1}{|\Lambda|}|\Lambda|\sum_{i\in Q_r}d_i=g_r,\quad


r=1,\dotsc,q.


\end{equation}


Таким образом, $\bold u\in X_4$. Теперь из ограниченности множества


$X_3$


следует, что множество $X_4$ непусто и ограничено.





Пусть $i\in Q_r$, $j\in P_k$. Выполняется равенство


$c^i_{\lambda(j)}=c^{\pi(i)}_j$ $\;\forall\lambda\in\Lambda$, где


$\pi$ ---


образ подстановки $\lambda$ в группе $\Pi$. Так как


$\{\pi(i)\mid\forall i\in Q_r\}=Q_r$, то


\begin{equation}


\begin{split}


\sum_{i\in Q_r}c^i_j&=\sum_{i\in Q_r}c^i_s \quad\forall j,s\in P_k;\\


\sum_{i\in Q_r}c^i_{j(k)}&=\frac{1}{|P_k|}\sum_{j\in P_k}\;


\sum_{i\in Q_r}c^i_j.


\end{split}


\end{equation}





Пусть $\bold u=(u_1,\dotsc,u_n)\in X_4$. Положим


$w_k=|P_k|u_j$, если $j\in P_k$; $\bold w=(w_1,\dotsc,w_p)$. Так как


$\bold u$ --- неподвижный по группе $\Lambda$ вектор, то с учетом (9)


получим


\begin{align*}


\la\bold f^r,\bold w\ra&=\sum^p_{k=1}\biggl(\,\frac{1}{|P_k|}


\sum_{j\in P_k}\;\sum_{i\in Q_r}c^i_j\biggr)w_k=\\


&=\sum_{i\in Q_r}\;\sum^p_{k=1}\;\sum_{j\in P_k}c^i_ju_j=


\sum_{i\in Q_r}\la\bold c^i,\bold u\ra=


|Q_r|\la\bold c^s,\bold u\ra\le g_r,\quad s\in Q_r,\quad


r=1,\dotsc,q.


\end{align*}


Следовательно, $\bold w\in X_2$.





Пусть $\bold w\in X_2$. Вектору $\bold w$ поставим в соответствие вектор


$\bold u=(u_1,\dotsc,u_n)$ такой, что


$|P_k|u_j=w_k$, если $j\in P_k$, $k=1,\dotsc,p$. Вектор $\bold u$


удовлетворяет условиям (7). Для всех $i\in Q_r$ имеем


\begin{align*}


|Q_r|\la\bold c^i,\bold u\ra&=\sum_{s\in Q_r}\la\bold c^s,\bold u\ra=


\sum_{s\in Q_r}\;\sum^p_{k=1}\;\sum_{j\in P_k}c^s_j


\frac{w_k}{|P_k|}=\\


&=\sum^p_{k=1}\biggl(\,\frac{1}{|P_k|}\sum_{s\in Q_r}\;


\sum_{j\in P_k}c^s_j\biggr)w_k=\\


&=\sum^p_{k=1}\biggl(\,\sum_{i\in Q_r}c^i_{j(k)}\biggr)w_k=


\la\bold f^r,\bold w\ra\le g_r,\quad r=1,\dotsc,q.


\end{align*}


Следовательно, $\bold u\in X_4$. Из установленного взаимно однозначного


отображения


множества $X_2$ на множество $X_4$ следует, что множество


$X_2$ непусто и


ограничено.


\end{pf}





\begin{thm}


Пусть $\bold v=(v_1,\dotsc,v_n)$ --- решение системы {\em (2), (3)},


$w_k=\sum\limits_{j\in P_k}v_j$.


Тогда $\bold w=(w_1,\dotsc,w_p)$ --- это решение системы {\em (4), (5)}.


\end{thm}





\begin{pf}


Из (2) следует (5). Так как $\bold v\in X_1$, то согласно формуле (8)


\begin{equation}


\sum_{\lambda\in\Lambda}\;\sum_{i\in Q_r}


\la\bold c^i,\bold v\lambda\ra\le |\Lambda|g_r,\quad


r=1,\dotsc,q.


\end{equation}


С другой стороны,


\begin{multline}


\sum_{\lambda\in\Lambda}\;\sum_{i\in Q_r}


\la\bold c^i,\bold v\lambda\ra=\sum_{i\in Q_r}\Bigl\la


\bold c^i,\sum_{\lambda\in\Lambda}\bold v\lambda\Bigr\ra=\\


=\sum_{i\in Q_r}\;\sum^n_{j=1}c^i_j\biggl(\,


\sum_{\lambda\in\Lambda}v_{\lambda(j)}\biggr)=\sum_{i\in Q_r}\;


\sum^p_{k=1}\;\sum_{j\in P_k}c^i_j


\frac{|\Lambda|}{|P_k|}\sum_{j\in P_k}v_{\lambda(j)}=\\


=|\Lambda|\sum^p_{k=1}\biggl(\,\frac{1}{|P_k|}


\sum_{j\in P_k}\;\sum_{i\in Q_r}c^i_j\biggr)w_k=


|\Lambda|\la\bold f^r,\bold w\ra,\quad r=1,\dotsc,q.


\end{multline}


Из (10), (11) следует, что условия (4) выполняются.


\end{pf}





Совместность системы (4), (5) является необходимым условием совместности


системы (2), (3). Пусть $X$, $Y$ --- множества решений систем (2), (3) и


(4), (5) соответственно. Согласно теореме~1


$$


X=\bigcup_{\bold y\in Y}X(\bold y),


$$


где


$X(\bold y)=\{\bold x\in X\bigl|\sum\limits_{j\in P_k}\bigr.x_j=y_k,


\;k=1,\dotsc,p\}$.





Дополнительные условия на переменные $x_j$, записываемые в виде линейных


уравнений, позволяют упростить решение системы (2), (3). Например, если в


системе (3) есть неравенства $x_j\ge 0$, $j=1,\dotsc,n$, и если $y_k=0$


для некоторого $\bold y=(y_1,\dotsc,y_p)\in Y$, то $x_j=0$ для


$\bold x=(x_1,\dotsc,x_n)\in X(\bold y)$, $j\in P_k$.





Если группа $\Lambda^0$ содержит две различные собственные подгруппы


$\Psi$, $\Phi$, то


может оказаться целесообразным следующий метод расчета. Пусть


$V$, $W$ ---


множества решений системы (4), (5) при $\Lambda=\Psi$, $\Lambda=\Phi$


соответственно; $G_k$, $k=1,\dotsc,s$, ---


орбиты группы $\Psi$ в множестве $\{1,\dotsc,n\}$;


$H_k$, $k=1,\dotsc,t$, --- орбиты


группы $\Phi$;


$X(\bold v,\bold w)$ --- множество решений системы (2), (3) при


дополнительных условиях


\begin{equation}


\sum_{j\in G_k}x_j=v_k,\quad k=1,\dotsc,s;\quad


\sum_{j\in H_k}x_j=w_k,\quad k=1,\dotsc,t.


\end{equation}


В соответствии с теоремой 1


\begin{equation}


X=\bigcup_{\bold v\in V}\;\bigcup_{\bold w\in W}


X(\bold v,\bold w).


\end{equation}





Возникает вопрос о ранге матрицы системы линейных уравнений (12).





\begin{thm}


Если $\Psi\Phi$ --- группа, то ранг матрицы системы {\em (12)} равен


$s+t-p_1$,


где $p_1$ --- число орбит группы $\Psi\Phi$.


\end{thm}





\begin{pf}


Группы $\Psi$, $\Phi$ --- это подгруппы группы $\Psi\Phi$. Так как любая


орбита


группы является объединением некоторых орбит подгруппы, то система (12)


распадается на $p_1$ независимых систем уравнений. Выделим одну из таких


систем, соответствующую некоторой орбите $P$ группы $\Psi\Phi$ и запишем


ее в виде


$\la\bold e^k,\bold x\ra=v_k$,


$k\in I=\{i\mid G_i\subseteq P\}$;


$\la\bold h^k,\bold x\ra=w_k$,


$k\in J=\{i\mid H_i\subseteq P\}$. Орбиты $G_k$, так же, как и орбиты


$H_k$, образуют разбиение орбиты $P$. Поэтому


$$


\sum_{k\in I}\bold e^k=\sum_{k\in J}\bold h^k.


$$


Следовательно, ранг матрицы системы, соответствующей орбите $P$, не


превышает $|I|+|J|-1$.





Предположим, что для некоторого $i\in I$ существуют такие числа


$\alpha_k$, $\beta_k$, что


$\alpha_i=0$,


$$


\sum_{k\in I}\alpha_k\bold e^k=\sum_{k\in J}\beta_k\bold h^k=\bold z.


$$


Так как векторы $\bold e^k$ неподвижны по группе $\Psi$, $\bold h^k$ ---


по группе $\Phi$, а


$\Psi\Phi$ --- группа, то вектор $\bold z$ неподвижен по группе


$\Psi\Phi$. Это значит, что


его компоненты $z_j$, $j\in P$, равны некоторому числу $\gamma$. Если


$\gamma\ne 0$,


то $\alpha_k\ne 0$ для


некоторого $k\in I$. Но поскольку $\alpha_i=0$, то вектор $\bold z$ не


является неподвижным по


группе $\Psi\Phi$. Из этого противоречия следует, что


$\gamma=0$; $\alpha_k=0$, $k\in I$; $\beta_k=0$, $k\in J$.


Таким


образом, ранг матрицы системы, соответствующей орбите $P$, не меньше


$|I|+|J|-1$,


следовательно, он равен $|I|+|J|-1$. Ранг матрицы системы (12) равен


сумме рангов матриц независимых систем. Эта сумма равна $s+t-p_1$.


\end{pf}





По поводу теоремы 2 заметим, что $\Psi\Phi$ --- группа, если и только


если $\Psi\Phi=\Phi\Psi$


(\cite{Kar}, с.23).





Нет смысла применять формулу (13), если одна из групп $\Psi$, $\Phi$


является подгруппой другой, т.к. в этом случае либо $s=p_1$, либо


$t=p_1$.





Представляет интерес определение условий, при которых система (4), (5)


оказывается симметричной.





\begin{thm}


Если $\Lambda$ --- нормальная собственная подгруппа группы


$\Lambda^0$ и если не существует группы $\Lambda^1$,


удовлетворяющей условиям:


{\em 1)}~$\Lambda<\Lambda^1\le\Lambda^0$, {\em 2)} орбиты групп


$\Lambda$, $\Lambda^1$


совпадают, то множество $F=\{\bold f^1,\dotsc,\bold f^q\}$ симметрично


по некоторой группе $\Psi^0$ подстановок, изоморфной фактор-группе


$\Lambda^0/\Lambda$.


\end{thm}





\begin{pf}


Пусть орбита $P^0$ группы $\Lambda^0$ является объединением орбит


$P_k$, $k\in I\subseteq\{1,\dotsc,p\}$,


группы $\Lambda$.


Действие фактор-группы $\Lambda^0/\Lambda$


переставляет равномощные орбиты $P_k$, $k\in I$, внутри


орбиты $P^0$. Действительно, т.к.


$\Lambda(P_k)=\{\lambda(j)\mid\,\forall\lambda\in\Lambda,\;


\forall j\in P_k\}=P_k$ и подгруппа


$\Lambda$ нормальна, то


$\mu(\Lambda(P_k))=\Lambda(\mu(P_k))=\mu(P_k)$ $\forall\mu\in\Lambda^0$.


Последнее равенство означает, что


$\mu(P_k)=\{\mu(j)\mid\;\forall j\in P_k\}$ --- это орбита группы


$\Lambda$, $\mu(P_k)\subseteq P^0$,


$|\mu(P_k)|=|P_k|$. Отсюда следует, что фактор-группа


$\Lambda^0/\Lambda$ порождает


некоторую группу


$\Psi^0$ подстановок на множестве $\{1,\dotsc,p\}$ номеров орбит группы


$\Lambda$. Элементу $\mu\Lambda\in\Lambda^0/\Lambda$


соответствует подстановка $\psi\in\Psi^0$ такая,


что $\mu\Lambda(P_k)=P_{\psi(k)}$.





Положим $\Lambda^1=\{\lambda\in\Lambda^0\mid\lambda(P_k)=P_k,\;


k=1,\dotsc,p\}$.


Непосредственно проверяется, что


$\Lambda^1$ --- группа, $\Lambda\le\Lambda^1\le\Lambda^0$. Группы


$\Psi^0$, $\Lambda^0/\Lambda$


изоморфны, если и только если


$\Lambda^1=\Lambda$. Следовательно, группы


$\Psi^0$, $\Lambda^0/\Lambda$


изоморфны, если не существует группы $\Lambda^1$,


удовлетворяющей условиям: 1)~$\Lambda<\Lambda^1\le\Lambda^0$,


2) орбиты групп $\Lambda$, $\Lambda^1$ совпадают.





Осталось доказать, что множество $F$ симметрично по группе $\Psi^0$.


Пусть $\psi\in\Psi^0$;


$\mu\Lambda$ --- прообраз подстановки $\psi$ в фактор-группе


$\Lambda^0/\Lambda$. Учитывая


формулу (9), получим


\begin{align}


f^r_{\psi(k)}&=\frac{1}{|P_{\psi(k)}|}\sum_{j\in P_{\psi(k)}}\;


\sum_{i\in Q_r}c^i_j=\notag\\


&=\frac{1}{|P_k|}\sum_{j\in P_k}\;\sum_{i\in Q_r}c^i_{\mu(j)}


=\frac{1}{|P_k|}\sum_{j\in P_k}\;\sum_{i\in Q_r}c^{\nu(i)}_j,


\end{align}


где $\nu=(\nu(1),\dotsc,\nu(m))$ --- образ подстановки


$\mu\in\Lambda^0$ в группе $\Pi^0$.


Согласно лемме 1 группа $\Lambda$ и группа $\Pi$, порождаемая группой


$\Lambda$ и


множеством $C$, изоморфны.


Следовательно, подгруппа $\Pi$ нормальна в $\Pi^0$, а фактор-группы


$\Pi^0/\Pi$, $\Lambda^0/\Lambda$


изоморфны. Действие фактор-группы $\Pi^0/\Pi$ переставляет орбиты $Q_r$.


Фактор-группа


$\Pi^0/\Pi$ порождает некоторую группу $\Phi^0$ подстановок на множестве


$\{1,\dotsc,q\}$ номеров


орбит группы $\Pi$. Элементу $\nu\Pi\in\Pi^0/\Pi$ соответствует


подстановка $\varphi\in\Phi^0$ такая, что


$\nu\Pi(Q_r)=\nu(Q_r)=Q_{\varphi(r)}$. Поэтому


\begin{equation}


\sum_{i\in Q_r}c^{\nu(i)}_j=\sum_{i\in Q_{\varphi(r)}}c^i_j.


\end{equation}


Из (14), (15) следует


$f^r_{\psi(k)}=f^{\varphi(r)}_k$,


$\bold f^r\psi=\bold f^{\varphi(r)}$, $r=1,\dotsc,q$. Таким образом,


множество $F$ симметрично


по группе $\Psi^0$, а группа $\Phi^0$, порождаемая группой $\Psi^0$ и


множеством $F$, изоморфна группе $\Psi^0$ согласно леммам 1, 2.


\end{pf}





Для произвольной группы $\Psi$ подстановок на множестве $\{1,\dotsc,p\}$


введем бинарное


отношение на пространстве $R^p$ равенством


$\bold y^1=\bold y^2\psi$ для некоторой подстановки


$\psi\in\Psi$.


Это отношение является отношением эквивалентности. Пусть $Y(\Psi)$ ---


множество


представителей всех классов эквивалентности множества $Y$ решений


системы (4), (5) относительно группы $\Psi$.





\begin{thm}


Если выполняются условия теоремы {\em 3} и если


$\bold d=(d_1,\dotsc,d_m)$ --- неподвижный по


группе $\Pi^0$ вектор, то


\begin{equation}


X=\biggl(\,\bigcup_{\bold y\in Y(\Psi^0)}X(\bold y)\biggr)\theta,


\end{equation}


где $\theta$ --- множество представителей смежных классов группы


$\Lambda^0$ по


подгруппе $\Lambda$.


\end{thm}





\begin{pf}


Если $\bold v\in X$, $\lambda\in\Lambda^0$, $i\in\{1,\dotsc,m\}$, то


$\la\bold c^i,\bold v\lambda\ra=\la\bold c^i\lambda^{-1},\bold v\ra=


\la\bold c^s,\bold v\ra\le d_s$,


где $i$, $s$ принадлежат одной и той же орбите


группы $\Pi^0$. По условию теоремы $d_i=d_s$, поэтому


$\la\bold c^i,\bold v\lambda\ra\le d_i$ и, следовательно,


$\bold v\lambda\in X$. Это значит, что множество $X$ симметрично по


группе $\Lambda^0$.





При $\varphi\in\Phi^0$, $r\in\{1,\dotsc,q\}$ орбиты $Q_r$,


$Q_{\varphi(r)}$ группы $\Pi$


содержатся в одной и той же орбите


группы $\Pi^0$. Поэтому $g_r=g_{\varphi(r)}$ и, следовательно,


$\bold g=(g_1,\dotsc,g_q)$ ---


неподвижный по группе $\Phi^0$


вектор. Теперь симметрия множества $Y$ по группе $\Psi^0$


устанавливается точно так же, как и симметрия множества


$X$ по группе $\Lambda^0$.





Пусть $\bold v\in X(w)$, $\bold w\in Y$, $\psi\in\Psi^0$;


$\mu\Lambda$ --- прообраз


подстановки $\psi$ в фактор-группе


$\Lambda^0/\Lambda$. По определению множества $X(y)$


\begin{equation}


\sum_{j\in P_k}v_{\mu(j)}=\sum_{j\in P_{\psi(k)}}v_j=


w_{\psi(k)},\quad k=1,\dotsc,p.


\end{equation}


Так как $\bold v\mu\in X$, $\bold w\psi\in Y$, то из (17) следует


$\bold v\mu\in X(\bold w\psi)$. По


определению множества $Y(\Psi^0)$


существует такая подстановка $\psi_0\in\Psi^0$, что


$\bold w\psi_0\in Y(\Psi^0)$. Тогда


$\bold v\mu_0\in X(\bold w\psi_0)$, где $\mu_0\Lambda$


--- прообраз


подстановки $\psi_0$ в фактор-группе


$\Lambda^0/\Lambda$, $\mu_0\in\Lambda^0$. Существует такая подстановка


$\nu\in\theta$, что


$\Lambda\mu^{-1}_0=\Lambda\nu$. Из определения множества


$X(\bold y)$, $\bold y\in Y$, следует,


что это множество


симметрично по группе $\Lambda$. Поэтому


$\bold v\in (X(\bold w\psi_0))\mu^{-1}_0=(X(\bold w\psi_0))\nu$. Это


значит, что любой вектор $\bold v\in X$


может быть определен по формуле (16).


\end{pf}





Симметричные по группе


$\Lambda^0$ системы (2) и (2), (3) назовем строго


симметричными, если $\bold d$ --- неподвижный по группе $\Pi^0$ вектор.


Некоторые


прикладные задачи формализуются симметричными или строго симметричными


системами линейных диофантовых неравенств, не требующими предварительной


симметризации. Такая ситуация имеет место в рассматриваемом ниже примере.





Задача определения гамильтонова цикла в неориентированном графе


$(V,E)$, где


$V=\{1,\dotsc,n\}$, формализуется следующей системой \cite{Vya}:


\begin{gather}


\sum^n_{j=1}x^i_j=1,\quad i=1,\dotsc,n;\quad


\sum^n_{i=1}x^i_j=1,\quad j=1,\dotsc,n;\\


\la\bold a^p,\bold x^i\ra+\sum_{k\in N_i}


\la\bold b^p,\bold x^k\ra\le 1,\quad i=1,\dotsc,n,\quad


p=0,1,\dotsc,n-1;\\


x^i_j\in\{0,1\},\quad i=1,\dotsc,n,\quad j=1,\dotsc,n.


\end{gather}


Здесь


$N_i=\{j\in V\mid\{i,j\}\notin E,\;i\ne j\}$,


$\bold x^i=(x^i_1,\dotsc,x^i_n)$, $\bold a^0=(1,1,0,\dotsc,0,1)$,


$\bold b^0=(1,0,\dotsc,0)$.


Векторы $\bold a^p$, $\bold b^p$ образованы циклическим


сдвигом на $p$ позиций вправо компонент векторов


$\bold a^0$, $\bold b^0$ соответственно.





Заметим, что систему (18)--(20) можно записать в точном соответствии с


записью системы (2), (3).





Система (18)--(20) строго симметрична по группе $\Omega\Delta_1$, где


$\Delta_1$ ---


циклическая группа порядка $n$ с порождающим элементом


$(n,1,2,\dotsc,n-1)$, $\Omega=\{(1,\dotsc,n),(n,n-1,\dotsc,2,1)\}$.


Действие группы $\Omega\Delta_1$ на $R^{n^2}$ определено следующим


образом:


$\bold x\lambda=(x^1_{\lambda(1)},\dotsc,x^1_{\lambda(n)},\dotsc,


x^n_{\lambda(1)},\dotsc,x^n_{\lambda(n)})$,


где $\bold x=(x^1_1,\dotsc,x^1_n,\dotsc,x^n_1,\dotsc,x^n_n)$,


$\lambda=(\lambda(1),\dotsc,\lambda(n))\in\Omega\Delta_1$.


Запишем систему (4), (5) для случая, когда $n$ --- четное число,


$\Lambda=\Delta_2$ ---


циклическая подгруппа порядка $\frac{n}{2}$ группы $\Delta_1$:


\begin{gather}


\sum^n_{i=1}y^i_1=\frac{n}{2};\quad


\sum^n_{i=1}y^i_2=\frac{n}{2};\\


2y^i_1+y^i_2+\sum_{k\in N_i}y^k_2\le\frac{n}{2};\quad


2y^i_2+y^i_1+\sum_{k\in N_i}y^k_1\le\frac{n}{2};\\


y^i_1+y^i_2=1;\quad y^i_1,y^i_2\in\{0,1\};\quad


i=1,\dotsc,n.


\end{gather}





Подгруппа $\Delta_2$ нормальна в группе $\Delta_1$. В соответствии с


теоремами 3, 4 (при


$\Lambda^0=\Delta_1$, $\Lambda=\Delta_2$)


система (21)--(23) строго симметрична по группе подстановок


порядка 2.





Систему (21)--(23) можно упростить, полагая


$M_i=\{j\mid\{i,j\}\in E\}$. С учетом того, что $|M_i\cup N_i|=n-1$,


получим эквивалентную систему:


\begin{gather}


\sum^n_{i=1}y^i_1=\frac{n}{2};\quad 2\le 3y^i_1+


\sum_{k\in M_i}y^k_1\le|M_i|+1;\\


y^i_1\in\{0,1\};\quad y^i_2=1-y^i_1;\quad


i=1,\dotsc,n.


\end{gather}


Любое решение системы (24), (25) содержит $n$ нулевых компонент.


Поэтому


согласно теореме 1 при расчетах по формуле (16) число переменных в


системе (18)--(20) сокращается в два раза.





Рассмотрим случай, когда $n$ делится на 3,


$\Lambda=\Omega\Delta_3$, $\Delta_3$ --- циклическая


подгруппа порядка $\frac{n}{3}$ группы $\Delta_1$. Пусть


$P^1_k$, $P^2_k$, $k=1,\dotsc,n$, ---


орбиты группы $\Omega\Delta_3$. В


данном случае орбите $P^1_k$ соответствуют переменные


$x^k_j$, $j=2+3i$, $i=0,1,\dotsc,\frac{n-3}{3}$; орбите


$P^2_k$


соответствуют остальные переменные


$x^k_j$; $|P^2_k|=2|P^1_k|=\frac{2n}{3}$. Система (4), (5) запишется в


следующем виде:


\begin{gather*}


\sum^n_{i=1}y^i_1=\frac{n}{3};\quad \sum^n_{i-1}y^i_2=\frac{2n}{3};\\


y^i_2+\sum_{k\in N_i}y^k_1\le\frac{n}{3};\quad


y^i_2+2y^i_1+\sum_{k\in N_i}y^k_2\le\frac{2n}{3};\\


y^i_1+y^i_2=1;\quad y^i_1,y^i_2\in\{0,1\};\quad i=1,\dotsc,n.


\end{gather*}


Упрощая эту систему, получим


\begin{gather*}


\sum^n_{i=1}y^i_1=\frac{n}{3};\quad 1\le 2y^i_1+


\sum_{k\in M_i}y^k_1\le|M_i|;\\


y^i_1\in\{0,1\};\quad y^i_2=1-y^i_1;\quad i=1,\dotsc,n.


\end{gather*}


Для подгрупп $\Delta_2$, $\Omega\Delta_3$ группы


$\Omega\Delta_1$ выполняется условие


теоремы 2, т.к.


$\Delta_2(\Omega\Delta_3)=(\Omega\Delta_3)\Delta_2=\Omega\Delta_1$.





\begin{thebibliography}{9}





\bibitem{Er}


Еремин И.И., Астафьев Н.Н. {\em Введение в теорию линейного и выпуклого


программирования}. Учеб. пособие. -- М.: Наука, 1976. -- 192 с.


\bibitem{Kar}


Каргаполов М.И., Мерзляков Ю.И. {\em Основы теории групп}. -- 3-е изд. --


М.: Наука, 1982. -- 288~с.


\bibitem{Vya}


Вьялицин А.А. {\em О перечислении гамильтоновых циклов} // Дискретная


математика. -- 1992. -- Т.3. -- Вып.3. -- С.46--49.








\end{thebibliography}





\vskip 2cm





\post{\it Северо-Казахстанский университет}{\it Поступила}


\post{}{03.05.1995}





\end{document}


